% ============================================
% Johns Hopkins Letter Template
% Assistant Professor Cover Letter – University of Vermont
% ============================================

\documentclass[11pt,letterpaper]{letter}

\usepackage{graphicx}
\usepackage{eso-pic}
\usepackage{geometry}
\usepackage{microtype}
\usepackage[hidelinks]{hyperref}

% ----- Page Setup -----
\geometry{left=25mm,right=25mm,top=25mm,bottom=25mm}

% ----- Logo Placement (foreground, first page only) -----
\newcommand{\PutJHULogoFG}[1]{%
  \AddToShipoutPictureFG*{%
    \AtPageUpperLeft{%
      \hspace{10mm}%
      \raisebox{-38mm}{%
        \includegraphics[width=6cm]{#1}%
      }%
    }%
  }%
}

% ----- Sender Information -----
\signature{Decory Edwards}
\address{Department of Economics \\ Johns Hopkins University \\ Baltimore, MD 21218 \\ \texttt{dedwar65@jh.edu}}

\begin{document}
\PutJHULogoFG{Images/jhulogo.png}

\begin{letter}{Search Committee \\
Department of Economics \\
University of Vermont \\
Burlington, VT 05405 \\ USA}

\opening{Dear Members of the Search Committee,}

I am writing to express my interest in the tenure‐track Assistant Professor position in Macroeconomics in the Department of Economics at the University of Vermont, starting Fall 2026. I am a Ph.D. candidate in Economics at Johns Hopkins University, advised by Professors Christopher D. Carroll, Nicholas W. Papageorge, and Stelios Fourakis, and I expect to receive my degree in May 2026. My research fields are macroeconomics, wealth and income dynamics, and computational economics.

My research agenda is unified by a focus on understanding the determinants of wealth inequality and the mechanisms through which heterogeneity in returns to wealth shapes aggregate outcomes. In my job market paper, I estimate the degree of heterogeneity in individual portfolio returns using micro-data from the Survey of Consumer Finances and simulate a structural model in which households face idiosyncratic return risk. I show that allowing for persistent heterogeneity in returns substantially improves the model’s ability to match the empirical distribution of wealth, offering a tractable way to connect micro-level return dispersion to macroeconomic inequality.

In a second project, I use the Health and Retirement Study to study the relationship between interpersonal trust and portfolio performance. Building on the literature that finds a hump-shaped relationship between trust and income, I show that a similar relationship holds for individual returns to wealth measured in the HRS data, highlighting a behavioral channel through which social attitudes shape saving and investment outcomes. This work also provides new estimates of the persistent component of returns to net worth, which motivate the calibration of heterogeneity in my structural model.

The University of Vermont’s emphasis on rigorous macroeconomics, quantitative methods, and high-quality undergraduate teaching aligns closely with my research and teaching agenda. The Department’s commitment to areas such as financial economics, international macroeconomics, growth, monetary economics, and empirical methods is a strong match for my experience and interests. In particular, I am excited by the possibility of contributing to undergraduate and elective courses such as \textit{Macroeconomic Theory}, \textit{Money and Banking}, \textit{Financial Markets and Institutions}, \textit{Economics of Uncertainty, Risk, and Finance}, and electives in advanced time series/financial econometrics. My empirical and computational strengths equip me to support UVM’s mission, including the work of centers such as the Vermont Complex Systems Center and the Vermont Advanced Computing Center.

\textbf{Teaching.} My teaching experience centers on undergraduate macroeconomics and an upper-level macro policy seminar, where I have led weekly discussion sections and held regular office hours for classes ranging from 16 to 40 students. My approach is student-centered: I aim to help students “convince themselves” that they have moved from confusion to understanding by holding structured office-hour coaching that builds trust and confidence. At UVM, I am prepared to teach intermediate macroeconomics, upper-level electives in my specialization, and quantitative methods for economics majors. I would also welcome developing an elective such as \textit{Computational Macroeconomics and Policy} or \textit{Wealth Dynamics and Macro Risk}, emphasizing reproducible workflows and evidence-based decision-making.

I believe I would be a strong fit because my computational and empirical toolkit allows me to (i) produce robust, data-backed estimates of returns and distributional outcomes, (ii) build and validate models that speak to corporate, regulatory, and policy-relevant questions at the intersection of business economics, macroeconomics and strategy, and (iii) collaborate across teams to present insights in concise, actionable formats for diverse audiences. I am enthusiastic about joining the collegial and intellectually vibrant community at UVM’s Department of Economics, contributing both to research and to undergraduate mentoring and advising.

I would be honored to contribute to the University of Vermont’s teaching and research mission, and to mentor students pursuing quantitative careers in economics. Thank you for your consideration; I welcome the opportunity to discuss my fit with the Department at your convenience.

\closing{Sincerely,}

\end{letter}
\end{document}
